% ---------------------------------------------------------------
% TCC1 - DETECÇÃO DE INTENÇÃO IMPLÍCITA PARA ASSISTENTES DE VOZ
% ---------------------------------------------------------------
% Phelipe Pereira de Souza
% SENAC Santo Amaro - Bacharelado em Ciência da Computação
% Orientador: Prof. Thyago Conchado Quintas
% ---------------------------------------------------------------

\documentclass[
    12pt,                % tamanho da fonte
    openright,           % capítulos começam em página ímpar
    oneside,             % para impressão em apenas um lado do papel
    a4paper,             % tamanho do papel
    english,             % idioma adicional
    brazil               % idioma principal
]{abntex2}

% ---------------------------------------------------------------
% PACOTES
% ---------------------------------------------------------------
\usepackage{lmodern}                % Fonte Latin Modern
\usepackage[T1]{fontenc}            % Seleção de códigos de fonte
\usepackage[utf8]{inputenc}         % Codificação do documento
\usepackage{lastpage}               % Usado pela Ficha catalográfica
\usepackage{indentfirst}            % Indenta o primeiro parágrafo
\usepackage{color}                  % Controle de cores
\usepackage{graphicx}               % Inclusão de gráficos
\usepackage{microtype}              % Melhorias de justificação
\usepackage{float}                  % Para posicionamento de figuras
\usepackage{booktabs}               % Tabelas profissionais
\usepackage{multirow}               % Células mescladas em tabelas
\usepackage{longtable}              % Tabelas longas
\usepackage{amsmath,amssymb}        % Símbolos matemáticos
\usepackage{pgfgantt}               % Diagramas de Gantt
\usepackage[brazilian,hyperpageref]{backref}  % Páginas com citações
\usepackage[alf]{abntex2cite}       % Citações ABNT

% ---------------------------------------------------------------
% CONFIGURAÇÕES DO PDF
% ---------------------------------------------------------------
\makeatletter
\hypersetup{
    pdftitle={\@title},
    pdfauthor={\@author},
    pdfsubject={TCC1 - SENAC Santo Amaro},
    pdfcreator={LaTeX with abnTeX2},
    pdfkeywords={tcc}{senac}{santo amaro}{trabalho de conclusão},
    colorlinks=true,        % false: links em caixas; true: links coloridos
    linkcolor=blue,         % cor dos links internos
    citecolor=blue,         % cor dos links para bibliografia
    filecolor=magenta,      % cor dos links para arquivos
    urlcolor=blue,
    bookmarksdepth=4
}
\makeatother

% ---------------------------------------------------------------
% CONFIGURAÇÕES DE APARÊNCIA DO PDF FINAL
% ---------------------------------------------------------------
\definecolor{blue}{RGB}{41,5,195}

% ---------------------------------------------------------------
% INFORMAÇÕES DO DOCUMENTO
% ---------------------------------------------------------------
\titulo{Detecção de Intenção Implícita para Ativação de Assistentes de Voz em Ambientes Controlados}
\autor{Phelipe Pereira de Souza}
\local{São Paulo}
\data{2030}
\orientador{Prof. Thyago Conchado Quintas}

\instituicao{%
  SENAC -- Serviço Nacional de Aprendizagem Comercial
  \par
  Unidade Santo Amaro
  \par
  Bacharelado em Ciência da Computação
}

\tipotrabalho{Trabalho de Conclusão de Curso I}

\preambulo{Proposta de Trabalho de Conclusão de Curso apresentado ao SENAC Santo Amaro como requisito parcial para aprovação na disciplina de TCC1 do Bacharelado em Ciência da Computação.}

% ---------------------------------------------------------------
% COMPILA O ÍNDICE
% ---------------------------------------------------------------
\makeindex

% ---------------------------------------------------------------
% INÍCIO DO DOCUMENTO
% ---------------------------------------------------------------
\begin{document}

% Seleciona o idioma do documento
\selectlanguage{brazil}

% Retira espaço extra obsoleto entre as frases
\frenchspacing

% ---------------------------------------------------------------
% ELEMENTOS PRÉ-TEXTUAIS
% ---------------------------------------------------------------
\pretextual

% ---------------------------------------------------------------
% CAPA
% ---------------------------------------------------------------
\imprimircapa

% ---------------------------------------------------------------
% FOLHA DE ROSTO
% ---------------------------------------------------------------
\imprimirfolhaderosto

% ---------------------------------------------------------------
% RESUMO
% ---------------------------------------------------------------
\begin{resumo}
Este trabalho apresenta a proposta de desenvolvimento de um sistema de detecção de intenção implícita para ativação de assistentes virtuais por voz, eliminando a necessidade de palavras de ativação explícitas (\textit{wake-words}). A abordagem proposta opera no nível textual, utilizando transcrições de fala obtidas via reconhecimento automático de fala (ASR) como entrada para um classificador de intenção baseado em modelos de linguagem pré-treinados para o português brasileiro, especificamente o BERTimbau. O sistema visa classificar automaticamente se uma fala é direcionada a um assistente virtual ou se pertence a uma conversação comum, em cenários com variáveis parcialmente controladas de ruído, número de participantes e contexto de interação. A pesquisa fundamenta-se em trabalhos recentes sobre \textit{Device-Directed Speech Detection} (DDSD) e classificação de intenção em linguagem natural, buscando contribuir para a área com uma solução voltada ao português brasileiro, idioma ainda pouco explorado na literatura sobre o tema. O trabalho está organizado em três capítulos que apresentam a introdução ao tema, a revisão bibliográfica e o desenvolvimento com a solução proposta e o cronograma de trabalho, seguidos das referências bibliográficas.

\vspace{\onelineskip}

\noindent
\textbf{Palavras-chave}: detecção de intenção. assistentes virtuais. processamento de linguagem natural. aprendizado de máquina. português brasileiro.
\end{resumo}

% ---------------------------------------------------------------
% SUMÁRIO
% ---------------------------------------------------------------
\pdfbookmark[0]{\contentsname}{toc}
\tableofcontents*
\cleardoublepage

% ---------------------------------------------------------------
% ELEMENTOS TEXTUAIS
% ---------------------------------------------------------------
\textual

% ===============================================================
% CAPÍTULO 1 - INTRODUÇÃO
% ===============================================================
\chapter{Introdução}
\label{cap:introducao}


Assistentes virtuais baseados em voz, como Google Assistant, Amazon Alexa e Apple Siri, tornaram-se componentes cada vez mais presentes no cotidiano, sendo incorporados em \textit{smartphones}, alto-falantes inteligentes, \textit{wearables} e sistemas de automação residencial. O mercado global de assistentes virtuais baseado em voz foi avaliado em US\$ 7,35 bilhões em 2024 e é projetado para alcançar US\$ 33,74 bilhões até 2030, com uma taxa de crescimento anual composta de 26,5\%, evidenciando a escala e a relevância dessa tecnologia \cite{NMSC2025}.

Na maioria dos sistemas atuais, a interação com assistentes virtuais baseados em voz é iniciada por meio de palavras de ativação explícitas, conhecidas como \textit{wake-words} ou \textit{hotwords}, tais como ``Hey Siri'', ``Alexa'' ou ``Ok Google''. Essas expressões atuam como um mecanismo de gatilho que sinaliza ao sistema o desejo do usuário de iniciar uma interação. Do ponto de vista técnico, essa abordagem é eficaz, pois permite reduzir o processamento contínuo de áudio e minimizar ativações indevidas \cite{Garg2021}. No entanto, a dependência de \textit{wake-words} impõe limitações significativas à naturalidade da comunicação humano-máquina, tornando a interação menos fluida e intuitiva quando comparada ao diálogo humano.

Em interações humanas cotidianas, comandos e solicitações são frequentemente expressos de forma implícita, inseridos naturalmente no fluxo da conversação, sem a necessidade de uma palavra específica para iniciar a comunicação. Em contraste, os sistemas baseados em \textit{wake-words} são vulneráveis ao fenômeno de ativações falsas por palavras foneticamente semelhantes, conhecido como \textit{FakeWake} \cite{Chen2021}. Esse problema compromete tanto a privacidade do usuário quanto a qualidade da experiência, reforçando as limitações da abordagem atual para uma interação verdadeiramente natural e intuitiva.

Nesse cenário, surge a área de pesquisa denominada Detecção de Fala Direcionada ao Dispositivo, em inglês \textit{Device-Directed Speech Detection} (DDSD), que visa determinar automaticamente se uma fala é direcionada a um assistente virtual ou se faz parte de uma conversa entre humanos \cite{Wagner2024}. Essa abordagem possibilita interações mais naturais e fluidas, dispensando a necessidade de gatilhos verbais explícitos (\textit{wake-words}) e representando um avanço significativo rumo a uma comunicação humano-máquina mais intuitiva e semelhante ao diálogo cotidiano.

Este trabalho propõe investigar a detecção de intenção implícita em interações por voz por meio de uma abordagem baseada no nível textual. O sistema proposto opera com transcrições de fala obtidas via Reconhecimento Automático de Fala, em inglês \textit{Automatic Speech Recognition} (ASR), e aplica técnicas de Processamento de Linguagem Natural (PLN) e Aprendizado de Máquina, em inglês \textit{Machine Learning} (ML), para classificar automaticamente se uma fala é direcionada ao dispositivo ou faz parte de uma conversa incidental entre humanos. A pesquisa concentra-se especificamente no português brasileiro, buscando contribuir para a democratização e a adaptação cultural dessa tecnologia.

% ---------------------------------------------------------------
\section{Contexto}
\label{sec:contexto}

O desenvolvimento dos assistentes virtuais baseados em voz acompanhou os avanços significativos nas áreas de ASR e PLN. Modelos como o Whisper \cite{Radford2023} representam um marco no campo do ASR, possibilitando transcrições de áudio com alta precisão e robustez em múltiplos idiomas, inclusive em condições adversas. Paralelamente, a arquitetura Transformer, proposta por \citeonline{Vaswani2017}, revolucionou o PLN, servindo de base para os modernos \textit{Large Language Models} (LLMs) que permitem uma compreensão semântica cada vez mais sofisticada da linguagem humana.

Para evitar o processamento contínuo de áudio e minimizar ativações indesejadas, a maioria dos assistentes virtuais atuais adota a técnica de detecção de palavra-chave (\textit{keyword spotting} ou \textit{wake-word detection}). Nessa abordagem, o sistema permanece em modo de baixo consumo até identificar uma \textit{wake-word} específica, como ``Hey Siri'', ``Alexa'' ou ``Ok Google'' \cite{Garg2021}. Embora eficiente do ponto de vista computacional e de privacidade, essa estratégia impõe uma comunicação rígida e pouco natural, exigindo que o usuário se adapte ao padrão do sistema em vez de interagir de forma fluida e conversacional.

Em cenários de uso cotidiano, é comum que os usuários expressem comandos e solicitações de maneira contextual e implícita, sem a necessidade de uma palavra de ativação explícita. Por exemplo, em um ambiente doméstico, uma pessoa pode dizer ``apaga a luz'' ou ``aumenta o volume da TV'', dirigindo-se ao dispositivo sem preceder o comando com uma \textit{wake-word}. Identificar de forma automática essa intenção implícita constitui o principal desafio da área de DDSD.

Trabalhos recentes na área de DDSD têm explorado abordagens multimodais que combinam características acústicas, textuais e visuais para determinar se uma fala é direcionada ao dispositivo \cite{Wagner2024, Krishna2023, Nie2021}. Embora eficazes, a maioria dessas pesquisas concentra-se no idioma inglês e utiliza arquiteturas complexas que requerem múltiplas modalidades de entrada. Este trabalho propõe uma abordagem simplificada e mais acessível, baseada unicamente no nível textual, com foco no português brasileiro e no uso de modelos de linguagem pré-treinados para esse idioma.

% ---------------------------------------------------------------
\section{Justificativa}
\label{sec:justificativa}

A crescente presença de assistentes virtuais baseados em voz nos mais variados dispositivos computacionais evidencia a relevância de pesquisas voltadas à melhoria de interação humano-máquina. Apesar dos avanços significativos em ASR e PLN, a maioria dos sistemas atuais ainda depende de palavras de ativação explícitas (\textit{wake-word}), o que compromete a naturalidade e a fluidez da comunicação.

Do ponto de vista acadêmico, o estudo de métodos alternativos de ativação baseados na detecção de intenção implícita contribui para o avanço em três áreas convergentes: Processamento de Linguagem Natural, Reconhecimento de Fala e Interação Humano-Computador. A literatura recente demonstra que abordagens multimodais para DDSD têm alcançado resultados promissores, com reduções de até 61\% na taxa de \textit{Equal Error Rate} (EER) em comparação com modelos unimodais \cite{Wagner2024}. Entretanto, tais abordagens geralmente demandam alto poder computacional e grandes volumes de dados multimodais, o que limita sua aplicação em dispositivos com recursos restritos. 

Sob a perspectiva prática, a capacidade de identificar automaticamente quando uma fala é direcionada ao assistente virtual, sem a necessidade de \textit{wake-words}, possibilita o desenvolvimento de interações mais naturais e intuitivas. Essa melhoria é especialmente relevante em ambientes residenciais e profissionais, onde o uso repetitivo de comandos de ativação pode causar fadiga cognitiva e desconforto ao usuário. 

Adicionalmente, identifica-se uma lacuna importante na literatura: a grande maioria dos trabalhos sobre DDSD foi desenvolvida para o idioma inglês. O português brasileiro, apesar de ser um dos idiomas mais falados do mundo, ainda é sub-representado nessa área de pesquisa. A recente disponibilidade de modelos como o BERTimbau \cite{Souza2020} e de datasets como FLA-Dataset \cite{Rocha2024} representa um avanço, mas ainda há poucas investigações que exploram essas ferramentas no contexto de detecção de intenção implícita.

Dessa forma, o presente trabalho busca contribuir para o preenchimento dessa lacuna ao investigar uma abordagem simplificada de detecção de intenção implícita baseada exclusivamente no nível textual, utilizando modelos de linguagem pré-treinados para o português brasileiro. Tal proposta visa oferecer uma solução mais acessível e eficiente, que não dependa de múltiplas modalidades de dados, facilitando sua aplicação em dispositivos com menor capacidade computacional.

% ---------------------------------------------------------------
\section{Objetivo}
\label{sec:objetivo}

\subsection{Objetivo principal}

Desenvolver e avaliar um sistema de classificação capaz de detectar intenção implícita em interações por voz, identificando, a partir de transcrições textuais, falas direcionadas a um assistente virtual em ambientes controlados, sem depender exclusivamente de palavras de ativação explícitas (\textit{wake-words}).

\subsection{Objetivos secundários}

\begin{enumerate}
    \item Realizar levantamento bibliográfico sobre Detecção de Fala Direcionada ao Dispositivo, em inglês \textit{Device-Directed Speech Detection} (DDSD), classificação de intenção em linguagem natural e modelos de linguagem pré-treinados;
    \item Construir um \textit{dataset} de transcrições de fala em português brasileiro, contendo exemplos rotulados de falas direcionadas ao assistente virtual e falas de conversação comum;
    \item Implementar um \textit{pipeline} de classificação de intenção composto por um módulo de Reconhecimento Automático de Fala, em inglês \textit{Automatic Speech Recognition} (ASR), seguido de um classificador textual baseado no BERTimbau, com modelos \textit{baseline} comparativos (TF-IDF + Regressão Logística e SVM linear);
    \item Avaliar o desempenho do sistema utilizando métricas de classificação como precisão (\textit{precision}), revocação (\textit{recall}), acurácia (\textit{accuracy}) e F1-\textit{score}, com ênfase na minimização de falsos positivos;
    \item Realizar análise comparativa conceitual entre a abordagem proposta e sistemas tradicionais baseados em \textit{wake-words}, utilizando dados de desempenho reportados na literatura.
\end{enumerate}

% ---------------------------------------------------------------
\section{Hipótese de pesquisa}
\label{sec:hipotese}

Este trabalho parte da hipótese de que um modelo de classificação baseado em características linguísticas extraídas de transcrições de fala em português brasileiro é capaz de identificar quando uma fala está sendo direcionada a um assistente virtual, mesmo na ausência de \textit{wake-words} explícitas.

% A hipótese central pode ser formalizada da seguinte forma: um modelo de classificação textual baseado no BERTimbau, treinado com transcrições de fala rotuladas em português brasileiro, é capaz de alcançar um F1-score superior a 0,85 na tarefa de distinguir falas direcionadas ao assistente de falas pertencentes a conversações comuns em ambientes controlados, mantendo a taxa de falso positivo abaixo de 10\%.

% ===============================================================
% CAPÍTULO 2 - REVISÃO BIBLIOGRÁFICA
% ===============================================================
\chapter{Revisão Bibliográfica}
\label{cap:revisao}

A revisão bibliográfica apresenta o embasamento teórico que fundamenta este trabalho, organizada por temas relevantes ao problema de pesquisa. Inicialmente, são apresentados os conceitos gerais que contextualizam a proposta. Em seguida, descreve-se a metodologia utilizada para a seleção dos trabalhos consultados. Os fundamentos técnicos são então detalhados, seguidos pela análise dos trabalhos diretamente relacionados ao problema abordado.

% ---------------------------------------------------------------
\section{Referencial teórico}
\label{sec:referencial}

\subsection{Assistentes virtuais e interação por voz}

Assistentes virtuais são sistemas computacionais projetados para interpretar comandos em linguagem natural e executar tarefas automatizadas em resposta às solicitações dos usuários. Esses sistemas integram tecnologias de Reconhecimento Automático de Fala, em inglês \textit{Automatic Speech Recognition} (ASR), Compreensão de Linguagem Natural, em inglês \textit{Natural Language Understanding} (NLU) e Texto para Fala, em inglês \textit{Text-to-Speech} (TTS) para viabilizar interações fluidas entre humanos e dispositivos \cite{Jurafsky2024}.

O funcionamento típico de um assistente virtual baseado em voz segue um \textit{pipeline} composto por etapas sequenciais: (1) captura do sinal de áudio; (2) Detecção de Atividade de Voz, em inglês \textit{Voice Activity Detection} (VAD); (3) ASR; (4) NLU; e (5) geração da resposta ou execução da ação solicitada. Na maioria dos sistemas comerciais, as etapas de processamento são precedidas por um módulo de detecção de \textit{wake-word}, que atua como gatilho para iniciar o processamento mais custoso \cite{Garg2021}.

\subsection{Detecção de Atividade de Voz (VAD)}

A VAD é uma etapa fundamental no pipeline de assistentes virtuais baseados em voz. Sua principal função é identificar os segmentos do sinal de áudio que contêm fala humana, separando-os de ruídos de fundo, silêncios e sons não-verbais. Ao realizar essa filtragem, o VAD evita o processamento desnecessário de trechos irrelevantes pelos módulos subsequentes de ASR e NLU, reduzindo significativamente o consumo computacional e a latência do sistema.

Em ambientes reais, a tarefa de VAD torna-se desafiadora devido à presença de ruídos variados (música, conversas paralelas, sons domésticos) e diferentes níveis de volume da voz. Modelos tradicionais baseavam-se em características acústicas como energia do sinal e coeficientes cepstrais de frequência Mel, em inglês \textit{Mel-Frequency Cepstral Coefficients} (MFCCs), enquanto abordagens mais recentes utilizam redes neurais profundas para maior robustez.

\cite{Li2021} propuseram um modelo de VAD baseado em Redes Neurais Convolucionais com Codificador Auditivo Mascarado, em inglês \textit{Masked Auditory Encoder Convolutional Neural Network} (M-AECNN), demonstrando que características inspiradas na audição humana podem melhorar significativamente a precisão da detecção em ambientes ruidosos.

\subsection{Reconhecimento automático de fala (ASR)}

O ASR consiste na conversão de sinais acústicos da fala humana em texto escrito. Os avanços recentes nessa área foram fortemente impulsionados pela adoção de arquiteturas baseadas em Transformer \cite{Vaswani2017}, que substituíram os modelos baseados em Redes Neurais Recorrentes, em inglês \textit{Recurrent Neural Networks} (RNNs), dominantes em períodos anteriores.

Um marco importante foi o modelo Whisper, proposto por \citeonline{Radford2023}. Treinado com aproximadamente 680 mil horas de dados de áudio multilíngue sob supervisão fraca, o Whisper demonstrou alta robustez e generalização, alcançando desempenho competitivo em tarefas de transcrição, tradução de fala e identificação de idioma, inclusive para o português brasileiro.

Paralelamente, tem crescido o interesse em soluções de ASR otimizadas para execução em dispositivos de borda (\textit{edge computing}). \citeonline{Xu2023} propuseram otimizações para modelos Conformer visando dispositivos com recursos extremamente limitados. Complementarmente, \citeonline{Jeffries2024} apresentaram o Moonshine, uma arquitetura de codificador-decodificador projetada especificamente para transcrição em tempo real e comandos de voz com baixa latência. Da mesma forma, \citeonline{Bevilacqua2024} desenvolveram o Whispy, um sistema que adapta modelos Whisper para operação em fluxos de áudio contínuos em tempo real.

\subsection{Processamento de Linguagem Natural (PLN) e Compreensão de Linguagem Natural (NLU)}

O PLN abrange o conjunto de técnicas computacionais destinadas à análise, compreensão e geração de linguagem humana. No contexto de assistentes virtuais, a subárea de NLU desempenha papel central, sendo responsável por extrair o significado semântico e a intenção por trás das entradas do usuário \cite{Jurafsky2024}.

Uma das tarefas fundamentais dentro da NLU é a classificação de intenção, responsável por identificar a ação ou intenção subjacente expressa em um enunciado do usuário. As abordagens para essa tarefa evoluíram de sistemas baseados em regras manuais para métodos estatísticos e, mais recentemente, para modelos baseados em aprendizado profundo.

Segundo o levantamento realizado por \citeonline{Weld2022}, os modelos baseados em Transformers pré-treinados, como o BERT e suas variantes, têm apresentado resultados superiores em tarefas de detecção de intenção e preenchimento de \textit{slots} (slot filling), superando abordagens tradicionais baseadas em redes recorrentes e convolucionais. A capacidade de realizar \textit{fine-tuning} desses modelos com quantidades relativamente pequenas de dados rotulados tornou essa estratégia especialmente atrativa para domínios com recursos limitados, como o português brasileiro.

\subsection{Modelos baseados em Transformers}

A arquitetura Transformer, proposta por \citeonline{Vaswani2017}, representou um marco no PLN ao introduzir o mecanismo de atenção multi-cabeça, em inglês \textit{multi-head attention}, como alternativa às tradicionais redes neurais recorrentes, em inglês \textit{recurrent neural networks} (RNNs). Diferentemente das RNNs, o Transformer permite o processamento paralelo de toda a sequência, resultando em maior eficiência computacional durante o treinamento e melhor capacidade de capturar dependências de longo alcance.

A partir dessa arquitetura, \citeonline{Devlin2019} propuseram o modelo BERT (\textit{Bidirectional Encoder Representations from Transformers}). O BERT realiza um pré-treinamento bidirecional por meio de duas tarefas principais: \textit{Masked Language Modeling} (MLM), que consiste em prever tokens mascarados dentro de uma sentença, e \textit{Next Sentence Prediction} (NSP), que treina o modelo a identificar se duas sentenças são consecutivas. Após o pré-treinamento, o modelo pode ser facilmente adaptado, utilizando técnicas de \textit{fine-tuning}, para tarefas específicas de PLN, adicionando uma camada de classificação sobre a representação do token especial \texttt{[CLS]}.

No contexto do português brasileiro, \citeonline{Souza2020} desenvolveram o BERTimbau, uma versão do BERT pré-treinada especificamente para o idioma. O modelo foi treinado em um \textit{corpus} de aproximadamente 2,7 bilhões de tokens extraídos da Wikipédia em português e do \textit{corpus} brWaC. Os resultados demonstraram que o BERTimbau supera os modelos multilíngues em diversas tarefas de PLN para o português, incluindo reconhecimento de entidades nomeadas, análise de similaridade textual e classificação de sentenças.

\subsection{Wake-words e seus limites}

O mecanismo de \textit{wake-words} (ou palavras de ativação) representa a estratégia predominante para o acionamento de assistentes virtuais baseados em voz. Consiste em manter um detector de palavras-chave, em inglês \textit{keyword spotter}, continuamente ativo no dispositivo, processando o áudio ambiente com baixo consumo computacional até que uma expressão pré-definida seja detectada \cite{Garg2021}.

Embora eficiente do ponto de vista computacional e de privacidade, esse mecanismo apresenta limitações importantes. \citeonline{Chen2021} investigaram o fenômeno conhecido como \textit{FakeWake}, no qual palavras foneticamente semelhantes à \textit{wake-word} podem acionar inadvertidamente o assistente. Os autores desenvolveram um gerador automático de expressões ``difusas'' (\textit{fuzzy words}) e conseguiram criar 965 expressões capazes de ativar indevidamente oito dos principais assistentes de voz comerciais em inglês e chinês, destacando a fragilidade dessa abordagem.

Para mitigar as ativações falsas, foram propostas técnicas de Mitigação de Falsos Gatilhos, em inglês \textit{False Trigger Mitigation} (FTM). \citeonline{Garg2021} apresentaram uma arquitetura unificada baseada em Transformer com processamento em fluxo contínuo, capaz de suprimir até 95\% dos falsos gatilhos utilizando apenas um segundo de áudio após a ativação. Complementarmente, \citeonline{Dighe2021} propuseram uma abordagem baseada em transferência de conhecimento, treinando um modelo compacto de LSTM (\textit{Long Short-Term Memory}) para identificar a intenção do usuário diretamente a partir de características acústicas, sem depender da transcrição completa via ASR.

Essas limitações inerentes aos sistemas baseados em \textit{wake-words} reforçam a importância de abordagens alternativas, como a DDSD, que buscam promover interações mais naturais e fluidas entre usuários e dispositivos.

\subsection{Detecção de fala direcionada ao dispositivo (DDSD)}

A DDSD consiste em uma tarefa de classificação binária cujo objetivo é distinguir automaticamente se uma fala é intencionalmente direcionada a um assistente virtual ou se faz parte de uma conversa incidental entre humanos ou ruído de fundo \cite{Krishna2023, Wagner2024}. Essa abordagem surge como uma alternativa promissora ao uso tradicional de \textit{wake-words}, possibilitando interações mais naturais e fluidas entre usuários e dispositivos.

As pesquisas recentes na área de DDSD têm se concentrado principalmente em abordagens multimodais. \citeonline{Wagner2024} demonstraram que a fusão de informações acústicas, hipóteses geradas pelo decodigicador de ASR e representações textuais processadas por \textit{Large Language Models} (LLMs) pode reduzir a taxa de erro em até 61\% em comparação com modelos unimodais. Da mesma forma, \citeonline{Krishna2023} exploraram a combinação de características verbais com informações prosódicas não-verbais (como entonação, ritmo e duração de pausas), obtendo melhorias de até 8,5\% na taxa de falsa aceitação.

Em outra linha de pesquisa multimodal, \citeonline{Nie2021} propuseram um sistema que integra informações visuais (detecção facial e rastreamento de lábios) com características acústicas (MFCCs) para determinar se o áudio capturado corresponde a uma fala direcionada ao dispositivo. Essa abordagem é especialmente útil em robôs de diálogo e dispositivos equipados com câmeras.

% ---------------------------------------------------------------
\section{Metodologia da pesquisa bibliográfica}
\label{sec:metodologia_pesquisa}

A pesquisa bibliográfica foi conduzida de forma sistemática com o objetivo de identificar trabalhos científicos relevantes relacionados à detecção de intenção em linguagem natural, assistentes virtuais baseados em voz, detecção de fala direcionada ao dispositivo (DDSD) e técnicas de aprendizado de máquina aplicadas ao reconhecimento e classificação de comandos de voz.

Foram consultadas bases de dados acadêmicas amplamente utilizadas na área de Computação, incluindo Google Scholar, Semantic Scholar, arXiv, IEEE Xplore e ACM Digital Library. Essas plataformas foram selecionadas por concentrarem artigos científicos, pré-publicações e trabalhos revisados por pares com abrangência nas áreas de processamento de fala, linguagem natural e aprendizado de máquina.

As buscas foram realizadas utilizando combinações de palavras-chave em inglês, tais como: \textit{``voice assistant''}, \textit{``device-directed speech detection''}, \textit{``intent detection''}, \textit{``keyword spotting''}, \textit{``implicit intent''}, \textit{``wake-word''}, \textit{``false trigger mitigation''}, \textit{``natural language understanding''} e \textit{``voice command classification''}. Essas expressões foram combinadas utilizando operadores booleanos (AND, OR) com o objetivo de ampliar ou restringir os resultados conforme necessário.

Como critérios de inclusão, foram considerados: (a) trabalhos publicados entre os anos de 2021 e 2026; (b) disponibilidade nos idiomas inglês ou português; (c) relação direta com o problema de pesquisa abordado, priorizando-se artigos que apresentassem experimentos, métodos de classificação ou avaliações de sistemas de interação por voz. Como exceção, foram incluídos artigos seminais anteriores ao período definido, por representarem contribuições fundamentais para a área (como os trabalhos sobre a arquitetura Transformer \cite{Vaswani2017} e o BERT \cite{Devlin2019}).

Como critérios de exclusão, foram desconsiderados: (a) trabalhos publicados antes de 2021, exceto artigos seminais; (b) artigos sem relação direta com detecção de intenção, DDSD ou classificação de comandos de voz; (c) trabalhos duplicados entre as bases consultadas.

O processo de seleção foi realizado em duas etapas. Na primeira etapa (triagem), foram identificados 118 artigos potencialmente relevantes com base na leitura do título e do resumo. Na segunda etapa (afunilamento), os artigos foram analisados com maior profundidade, resultando na seleção de 18 trabalhos para compor a fundamentação teórica, os trabalhos relacionados e as referências deste estudo. Adicionalmente, foram incluídos 5 artigos seminais anteriores a 2021 pela sua importância para a base teórica do trabalho.

% ---------------------------------------------------------------
\section{Fundamentos}
\label{sec:fundamentos}

Esta seção detalha as tecnologias, algoritmos e metodologias específicas que serão empregados no desenvolvimento da solução proposta.

\subsection{Arquitetura Transformer}

A arquitetura Transformer \cite{Vaswani2017} é composta por um codificador (\textit{encoder}) e um decodificador (\textit{decoder}), ambos formados por camadas empilhadas que utilizam mecanismos de atenção multi-cabeça e redes neurais \textit{feed-forward}. O mecanismo de atenção calcula a relevância de cada elemento da sequência de entrada em relação aos demais, por meio das matrizes de consultas ($Q$), chaves ($K$) e valores ($V$):

\begin{equation}
\text{Attention}(Q, K, V) = \text{softmax}\left(\frac{QK^T}{\sqrt{d_k}}\right)V
\label{eq:attention}
\end{equation}

\noindent onde $d_k$ é a dimensão do vetor de chaves. A atenção multi-cabeça executa essa operação em paralelo com diferentes projeções lineares, permitindo que o modelo capture diferentes tipos de dependências simultaneamente.

Para a tarefa de classificação de intenção proposta neste trabalho, será utilizada apenas a parte codificadora (\textit{encoder}) da arquitetura, conforme o modelo BERT.

\subsection{BERT e BERTimbau}

O BERT \cite{Devlin2019} utiliza exclusivamente o codificador do Transformer para gerar representações contextuais bidirecionais de texto. O pré-treinamento é realizado em grandes \textit{corpora} de texto não rotulado, e o modelo resultante pode ser ajustado (\textit{fine-tuning}) para tarefas específicas com a adição de camadas de classificação.

Para a classificação de sequências (como a classificação de intenção), a representação do token especial \texttt{[CLS]}, posicionado no início de cada entrada, é utilizada como representação global da sentença. Essa representação é então passada a uma camada densa com ativação \textit{softmax} para produzir a predição final.

O BERTimbau \cite{Souza2020} é a versão do BERT pré-treinada especificamente para o português brasileiro. Disponível nas variantes \textit{base} (110 milhões de parâmetros) e \textit{large} (335 milhões de parâmetros), o modelo foi treinado com textos em português e demonstrou desempenho superior ao BERT multilíngue em tarefas de PLN para o português. Neste trabalho, será utilizada a variante \textit{base}, considerando as restrições de recursos computacionais disponíveis (GPU RTX A2000 com 12 GB de memória).

\subsection{Modelos baseline}

Para fins de comparação, serão implementados dois modelos \textit{baseline} que representam abordagens tradicionais de classificação de texto:

\begin{itemize}
    \item \textbf{TF-IDF + Regressão Logística}: a representação TF-IDF (\textit{Term Frequency-Inverse Document Frequency}) converte cada texto em um vetor numérico que pondera a importância de cada termo no documento em relação ao \textit{corpus}. Esse vetor é utilizado como entrada para um classificador de Regressão Logística;
    \item \textbf{TF-IDF + SVM Linear}: a mesma representação TF-IDF é utilizada como entrada para uma Máquina de Vetores de Suporte (\textit{Support Vector Machine}) com \textit{kernel} linear, que busca encontrar o hiperplano ótimo de separação entre as classes.
\end{itemize}

Esses modelos servem como referência para avaliar o ganho de desempenho proporcionado pelo uso de representações contextuais do BERTimbau em comparação com representações estáticas baseadas em frequência de termos.

\subsection{Métricas de avaliação}

O desempenho dos modelos será avaliado com as seguintes métricas:

\begin{itemize}
    \item \textbf{Acurácia} (\textit{accuracy}): proporção de classificações corretas em relação ao total de amostras;
    \item \textbf{Precisão} (\textit{precision}): proporção de verdadeiros positivos entre as amostras classificadas como positivas, indicando a confiabilidade das ativações do sistema;
    \item \textbf{Revocação} (\textit{recall}): proporção de verdadeiros positivos entre todas as amostras efetivamente positivas, indicando a capacidade do sistema de detectar comandos reais;
    \item \textbf{F1-\textit{score}}: média harmônica entre precisão e revocação, fornecendo uma métrica balanceada;
    \item \textbf{Matriz de confusão}: representação tabular dos resultados de classificação, permitindo a análise detalhada de falsos positivos e falsos negativos.
\end{itemize}

Neste trabalho, será dada ênfase especial à minimização de falsos positivos, uma vez que ativações indevidas do assistente virtual comprometem a privacidade do usuário e a confiabilidade do sistema.

\subsection{Reconhecimento automático de fala com Whisper}

O módulo de ASR do \textit{pipeline} proposto utilizará o modelo Whisper \cite{Radford2023} para a conversão do áudio em texto. O Whisper é um modelo \textit{encoder-decoder} baseado em Transformer, treinado com supervisão fraca em grande escala. Para este trabalho, será utilizada a variante \textit{small} ou \textit{medium} do modelo, que oferece suporte ao português e pode ser executada na GPU disponível.

O fluxo de processamento seguirá a sequência: captura do áudio $\rightarrow$ transcrição via Whisper $\rightarrow$ pré-processamento textual (normalização, remoção de ruído textual) $\rightarrow$ classificação de intenção via BERTimbau.

% ---------------------------------------------------------------
\section{Trabalhos relacionados}
\label{sec:trabalhos_relacionados}

Esta seção apresenta três trabalhos que compartilham o mesmo foco de resolução de problema: a detecção automática de fala direcionada a assistentes virtuais sem o uso de \textit{wake-words}. Para cada trabalho, são descritos o problema abordado, a metodologia utilizada, os principais resultados e as limitações identificadas.

\subsection{Wagner et al. (2024) --- Abordagem multimodal para DDSD com LLMs}

\citeonline{Wagner2024} abordaram o problema de determinar se uma fala é direcionada a um assistente virtual sem o uso de uma expressão de ativação prévia --- o mesmo problema central deste trabalho. Para resolvê-lo, os autores propuseram uma abordagem multimodal que combina três fontes de informação: (1) representações acústicas obtidas de um codificador de áudio; (2) hipóteses de transcrição do decodificador ASR; e (3) sinais do decodificador, todos integrados em um Large Language Model (LLM). O modelo foi treinado utilizando técnicas de adaptação de baixo posto (\textit{Low-Rank Adaptation} --- LoRA) e \textit{prefix tuning}.

Os resultados mostraram reduções na Taxa de Erro de Igualdade (EER) de até 39\% em relação a modelos baseados apenas em texto e de até 61\% em relação a modelos baseados apenas em áudio. Aumentar o tamanho do LLM e aplicar LoRA proporcionou reduções adicionais de até 18\% no EER. No entanto, a solução depende de um LLM com capacidade significativa de processamento e requer acesso simultâneo a múltiplas modalidades de dados, o que pode limitar sua aplicação em cenários com recursos restritos.

O presente trabalho diferencia-se dessa abordagem ao propor uma solução exclusivamente textual, baseada em um modelo de linguagem específico para o português brasileiro (BERTimbau), que não requer informações acústicas ou sinais do decodificador, sendo portanto mais acessível em termos de complexidade e recursos.

\subsection{Nie et al. (2021) --- Ativação multimodal sem wake-words}

\citeonline{Nie2021} abordaram o mesmo problema de ativar um robô de diálogo sem a necessidade de \textit{wake-words}. A solução proposta, denominada \textit{Multimodal Activation Scheme} (MAS), combina dois componentes: (1) detecção de consistência audiovisual, que verifica se o áudio captado corresponde aos movimentos labiais do usuário identificado pela câmera; e (2) inferência de intenção semântica de conversa, que analisa a transcrição da fala para determinar se o conteúdo é direcionado ao dispositivo.

Para a detecção audiovisual, foram utilizadas duas redes CNN heterogêneas para processar, respectivamente, características faciais (\textit{landmarks}) e características espectrais (MFCC). Para a inferência de intenção, a transcrição foi processada por fatoração de matrizes para identificar tópicos latentes de interação humano-robô. Os autores construíram um \textit{dataset} composto por 12.741 vídeos gravados por 194 voluntários.

A solução demonstrou eficácia na combinação de modalidades, porém apresenta limitações relacionadas à necessidade de câmera frontal e à complexidade do processamento visual em tempo real. Além disso, o sistema foi desenvolvido e avaliado exclusivamente em chinês, sem validação em outros idiomas. O presente trabalho diferencia-se ao focar exclusivamente na modalidade textual, dispensando sensores visuais, e ao direcionar a pesquisa para o português brasileiro.

\subsection{Wang et al. (2024) --- M$^{3}$V: abordagem multi-view para DDSD}

\citeonline{Wang2024} propuseram o M$^{3}$V (\textit{Multi-modal Multi-view}), uma abordagem que trata o problema de DDSD como uma tarefa de aprendizado multi-\textit{view}. O diferencial dessa abordagem é o tratamento explícito dos erros introduzidos pelo sistema ASR, que frequentemente causam desalinhamento entre as representações textuais e acústicas.

O modelo introduz visões unimodais (processamento individual de áudio e texto) e uma visão de alinhamento texto-áudio, além da visão multimodal tradicional. Os resultados experimentais mostraram que o M\textsuperscript{3}V supera significativamente modelos treinados com uma única modalidade e até mesmo supera o julgamento humano em dados com erros de ASR.

As limitações incluem a complexidade arquitetural do modelo e a dependência de dados multimodais pareados para treinamento. O presente trabalho diferencia-se ao adotar uma abordagem unimodal (apenas texto), que, embora potencialmente menos precisa, é substancialmente mais simples de implementar e não requer o tratamento explícito de desalinhamentos entre modalidades.

\vspace{0.5cm}

A Tabela~\ref{tab:trabalhos_relacionados} sintetiza a comparação entre os trabalhos relacionados e a proposta deste trabalho.

\begin{table}[htb]
    \centering
    \caption{Comparação entre trabalhos relacionados}
    \label{tab:trabalhos_relacionados}
    \begin{tabular}{l|c|c|c|c}
        \toprule
        \textbf{Critério} & \textbf{Wagner} & \textbf{Nie} & \textbf{Wang} & \textbf{Este trabalho} \\
        \midrule
        Modalidades  & Áudio + Texto + ASR & Áudio + Vídeo & Áudio + Texto & Texto \\
        \midrule
        Modelo       & LLM + LoRA    & CNN + MF      & Multi-view    & BERTimbau \\
        \midrule
        Idioma       & Inglês        & Chinês        & Chinês        & Português (BR) \\
        \midrule
        Dataset      & Interno       & 12.741 vídeos & Interno       & Construído (PT-BR) \\
        \midrule
        Validação    & EER           & Acurácia      & EER           & F1, Precision, Recall \\
        \midrule
        Limitações   & Requer LLM    & Requer câmera & Complexidade  & Unimodal \\
        \bottomrule
    \end{tabular}
    \legend{Fonte: Elaborado pelo autor (2026)}
\end{table}

% ===============================================================
% CAPÍTULO 3 - DESENVOLVIMENTO
% ===============================================================
\chapter{Desenvolvimento}
\label{cap:desenvolvimento}

Neste capítulo, apresenta-se a solução proposta para o problema identificado, bem como o cronograma de trabalho previsto para o desenvolvimento completo do projeto no TCC2.

% ---------------------------------------------------------------
\section{Solução proposta}
\label{sec:solucao}

A solução proposta consiste em um sistema de detecção de intenção implícita baseado em um \textit{pipeline} de processamento textual, conforme ilustrado a seguir:

\begin{center}
\textbf{Áudio} $\rightarrow$ \textbf{ASR (Whisper)} $\rightarrow$ \textbf{Texto} $\rightarrow$ \textbf{Pré-processamento} $\rightarrow$ \textbf{Classificador (BERTimbau)} $\rightarrow$ \textbf{Saída}
\end{center}

\subsection{Visão geral}

O sistema receberá como entrada segmentos de áudio capturados em ambientes controlados e produzirá como saída uma classificação binária indicando se a fala é \textbf{direcionada ao assistente virtual} ou \textbf{pertence a uma conversação comum}. O processamento é dividido em dois módulos principais:

\begin{itemize}
    \item \textbf{Módulo ASR}: responsável pela transcrição do áudio para texto, utilizando o modelo Whisper \cite{Radford2023} em sua variante adequada para o português brasileiro;
    \item \textbf{Módulo NLU/Classificador}: responsável pela análise do texto transcrito e classificação da intenção, utilizando o BERTimbau \cite{Souza2020} com \textit{fine-tuning} para a tarefa de classificação binária.
\end{itemize}

\subsection{Definição de ambiente controlado}

Para fins deste trabalho, o ambiente controlado é definido como um cenário com variáveis parcialmente controladas de ruído, número de participantes e contexto de interação, incluindo:

\begin{itemize}
    \item Ambientes residenciais simulados (sala de estar, quarto);
    \item Ambientes profissionais (escritório com dois ou mais ocupantes);
    \item Número limitado de falantes simultâneos (1 a 3 pessoas);
    \item Ausência de sobreposição intensa de fala;
    \item Ruído de fundo leve a moderado (televisão, ventilador, conversas ao fundo).
\end{itemize}

\subsection{Construção do dataset}

O \textit{dataset} será construído com foco exclusivo no português brasileiro, composto por transcrições de fala organizadas em duas classes:

\begin{enumerate}
    \item \textbf{Device-directed} (DD): transcrições de comandos e solicitações direcionados a um assistente virtual (exemplos: ``apaga a luz'', ``qual a previsão do tempo'', ``toca uma música'');
    \item \textbf{Non-device-directed} (NDD): transcrições de falas pertencentes a conversações comuns (exemplos: ``que horas a gente vai sair?'', ``eu preciso ir ao mercado'', ``gostei muito do filme'').
\end{enumerate}

A construção do \textit{dataset} seguirá três estratégias complementares: (a) tradução e adaptação de \textit{datasets} existentes em inglês, como o Fluent Speech Commands; (b) criação de exemplos originais baseados em cenários de uso residencial e profissional; e (c) validação e enriquecimento por meio da classificação de transcrições de vídeos no YouTube em português brasileiro.

O FLA-Dataset \cite{Rocha2024}, um banco de dados de comandos de áudio em português brasileiro baseados em localização, será também utilizado como referência para a construção de exemplos de comandos em português.

\subsection{Metodologia de desenvolvimento}

O desenvolvimento seguirá uma abordagem iterativa, organizada nas seguintes etapas:

\begin{enumerate}
    \item Revisão bibliográfica e definição dos requisitos;
    \item Construção e curadoria do \textit{dataset} em português brasileiro;
    \item Implementação do \textit{pipeline} ASR (Whisper) + pré-processamento;
    \item Treinamento e avaliação dos modelos \textit{baseline} (TF-IDF + Regressão Logística e TF-IDF + SVM);
    \item \textit{Fine-tuning} e avaliação do modelo BERTimbau;
    \item Análise comparativa dos resultados;
    \item Documentação e escrita do TCC2.
\end{enumerate}

\subsection{Critérios de validação}

O sistema será validado com base nas seguintes métricas: acurácia, precisão, revocação e F1-\textit{score}. A ênfase será dada à minimização de falsos positivos (ativações indevidas), métrica crítica para a usabilidade e confiabilidade de um assistente virtual.

Adicionalmente, será realizada uma análise comparativa conceitual com sistemas baseados em \textit{wake-words}, utilizando dados de desempenho reportados na literatura, como as taxas de falso positivo de sistemas comerciais e os resultados de mitigação de falsos gatilhos apresentados por \citeonline{Garg2021} e \citeonline{Chen2021}.

\subsection{Tecnologias e ferramentas}

\begin{itemize}
    \item \textbf{Linguagem}: Python 3.10+
    \item \textbf{ASR}: Whisper (OpenAI) --- variante \textit{small} ou \textit{medium}
    \item \textbf{Modelo principal}: BERTimbau \textit{base} \cite{Souza2020} via biblioteca Hugging Face Transformers
    \item \textbf{Modelos baseline}: scikit-learn (TF-IDF, Regressão Logística, SVM)
    \item \textbf{Treinamento}: PyTorch com aceleração GPU (RTX A2000, 12 GB)
    \item \textbf{Pré-processamento}: NLTK, spaCy (português)
    \item \textbf{Versionamento}: Git
\end{itemize}

% ---------------------------------------------------------------
\section{Cronograma de trabalho}
\label{sec:cronograma}

O cronograma a seguir apresenta, no formato de diagrama de Gantt, as atividades previstas para o desenvolvimento do TCC2, distribuídas ao longo das semanas do semestre.

\begin{figure}[htb]
    \centering
    \caption{Cronograma de atividades --- Diagrama de Gantt}
    \label{fig:gantt}
    \begin{ganttchart}[
        hgrid,
        vgrid,
        x unit=0.72cm,
        y unit title=0.6cm,
        y unit chart=0.7cm,
        title height=1,
        bar height=0.5,
        bar label font=\scriptsize,
        title label font=\small,
        milestone label font=\scriptsize\itshape,
        group label font=\scriptsize\bfseries,
        bar/.append style={fill=blue!40},
        group/.append style={draw=black, fill=blue!70},
        milestone/.append style={fill=red!70},
    ]{1}{16}
        % Título: semanas de 1 a 16
        \gantttitle{Semanas}{16} \\
        \gantttitlelist{1,...,16}{1} \\

        % Grupo 1: Dataset e Preparação
        \ganttgroup{Fase 1: Dados}{1}{5} \\
        \ganttbar{Construção dataset}{1}{4} \\
        \ganttbar{Pré-processamento}{3}{5} \\
        \ganttmilestone{Dataset pronto}{5} \\

        % Grupo 2: Desenvolvimento dos Modelos
        \ganttgroup{Fase 2: Modelagem}{5}{11} \\
        \ganttbar{Pipeline ASR (Whisper)}{5}{6} \\
        \ganttbar{Modelos baseline}{6}{8} \\
        \ganttbar{Fine-tuning BERTimbau}{7}{10} \\
        \ganttbar{Análise comparativa}{10}{11} \\
        \ganttmilestone{Modelos avaliados}{11} \\

        % Grupo 3: Validação e Escrita
        \ganttgroup{Fase 3: Validação}{12}{16} \\
        \ganttbar{Testes e validação}{12}{14} \\
        \ganttbar{Ajustes e otimização}{13}{15} \\
        \ganttbar{Redação final TCC2}{13}{16} \\
        \ganttmilestone{Entrega final}{16}

        % Ligações
        \ganttlink{elem2}{elem3}
        \ganttlink{elem3}{elem4}
        \ganttlink{elem6}{elem7}
        \ganttlink{elem7}{elem8}
        \ganttlink{elem8}{elem9}
        \ganttlink{elem9}{elem10}
        \ganttlink{elem12}{elem13}
    \end{ganttchart}
    \legend{Fonte: Elaborado pelo autor (2026)}
\end{figure}

% ---------------------------------------------------------------
% ELEMENTOS PÓS-TEXTUAIS
% ---------------------------------------------------------------
\postextual

% ===============================================================
% REFERÊNCIAS BIBLIOGRÁFICAS
% ===============================================================
\bibliography{bibliografia}

\end{document}
